
\subsubsection{Truthfulness Benchmarks}

TruthfulQA \citep{lin2022truthfulqa} introduced 817 questions across 38 categories designed to elicit imitative falsehoods---plausible-sounding misconceptions that models reproduce from training data. The benchmark specifically targets misconceptions that humans commonly believe, making it orthogonal to factual knowledge tests.

HaluEval \citep{li2023halueval} provides 35,000 samples across QA, dialogue, and summarization tasks, targeting hallucinated details during text generation. Unlike TruthfulQA's focus on misconceptions, HaluEval tests generation coherence and consistency maintenance.

FactScore \citep{min2023factscore} introduced atomic fact decomposition for evaluating long-form generation. By breaking responses into individual claims and verifying each against retrieval sources, FactScore provides fine-grained factuality measurement. This methodology differs from both multiple-choice evaluation and binary detection.

These benchmarks developed independently, each establishing its evaluation paradigm without cross-benchmark validation.

\subsubsection{Benchmark Meta-Evaluation}

BenchBench \citep{perlitz2024benchbench} proposed a meta-benchmark methodology for evaluating benchmark agreement. Their finding that benchmark agreement varies with methodological choices motivates investigation of truthfulness-specific correlation structure.

Ailem et al. (2024) demonstrated non-random correlations in model performance across test prompts within benchmarks, showing that accounting for prompt-level correlations can change model rankings. This work operates at the prompt level within benchmarks; the present study extends to model-level correlations across benchmarks.

\subsubsection{Position of This Work}

This work differs from prior studies in three ways: (1) domain-specific correlation analysis for truthfulness benchmarks on the same model population, (2) model-level cross-benchmark analysis rather than prompt-level within-benchmark correlations, and (3) explicit multi-dimensional structure testing using factor analysis.
